% =============================================================================
% Section 2: Data & Construction
% =============================================================================
\section{Data and Construction}\label{sec:data-construction}

This section provides a concise overview of the ISI's data inputs,
aggregation rule, and classification thresholds.  Full construction
details---including normalisation procedures, missing-data protocols,
fallback basis rules, and the integrity chain---are documented in
Paper~1~\cite{isi_paper1}.

\subsection{Scope and vintage}

The results in this paper are computed from the ISI v1.0 snapshot for
vintage year 2024.  The scope covers $N = 27$ EU member states.  All
27~countries have complete data across all six axes
(\texttt{complete = true}).

\subsection{Axes and data windows}

The ISI comprises six axes, each measuring counterparty concentration
in a distinct domain of cross-border exposure.  \Cref{tab:axis-sources}
summarises each axis, its primary data source, and the observation
window used for the 2024 vintage.

\begin{table}[htbp]
\centering
\caption{ISI axes: definitions, primary data sources, and observation
  windows for vintage~2024.}
\label{tab:axis-sources}
\small
\begin{tabular}{@{}clll@{}}
\toprule
\textbf{Axis} & \textbf{Domain} & \textbf{Primary Source}
  & \textbf{Window} \\
\midrule
1 & Financial sovereignty
  & BIS LBS, IMF CPIS~\cite{bis_lbs,imf_cpis}
  & 2022--2024 \\
2 & Energy dependency
  & IEA, Eurostat~\cite{iea_wes,eurostat_comext}
  & 2022--2024 \\
3 & Technology / semiconductor dependency
  & Eurostat Comext~\cite{eurostat_comext}
  & 2022--2024 \\
4 & Defence industrial dependency
  & SIPRI TIV~\cite{sipri_transfers}
  & 2019--2024 \\
5 & Critical inputs / raw materials
  & UN Comtrade, Eurostat~\cite{uncomtrade,eurostat_comext}
  & 2022--2024 \\
6 & Logistics / freight dependency
  & Eurostat Comext~\cite{eurostat_comext}
  & 2022--2024 \\
\bottomrule
\end{tabular}
\end{table}

\noindent\textit{Note.}  The defence axis uses a six-year rolling
window (2019--2024) rather than the three-year window applied to other
axes.  This reflects the lower frequency and higher lumpiness of
arms-transfer flows, as documented in Paper~1~\cite{isi_paper1}.

\subsection{Aggregation rule}

The ISI composite for country~$i$ is the unweighted arithmetic mean
of the six axis-level scores:
\begin{equation}\label{eq:composite}
  \text{ISI}_i
  = \frac{1}{6}\sum_{j=1}^{6} A_{j,i}
  \,,\qquad
  A_{j,i} \in [0,\,1]
  \,,
\end{equation}
where $A_{j,i}$ denotes the normalised score of country~$i$ on axis~$j$.
Each axis score is bounded to the unit interval by the normalisation
procedure defined in Paper~1.

The equal-weight specification is the baseline.  \Cref{sec:robustness}
examines sensitivity to alternative weighting schemes through Dirichlet
perturbation.

\subsection{Classification thresholds}\label{sec:classification}

Countries are classified into four tiers based on the composite score:

\begin{table}[htbp]
\centering
\caption{ISI classification tiers.}
\label{tab:classification-tiers}
\begin{tabular}{@{}lll@{}}
\toprule
\textbf{Tier} & \textbf{Composite range} & \textbf{Count (EU-27)} \\
\midrule
Highly concentrated     & $\text{ISI} \geq 0.50$ & 1 (Malta) \\
Moderately concentrated & $0.25 \leq \text{ISI} < 0.50$ & 23 \\
Mildly concentrated     & $0.15 \leq \text{ISI} < 0.25$ & 3 (Latvia, Slovakia, France) \\
Unconcentrated          & $\text{ISI} < 0.15$ & 0 \\
\bottomrule
\end{tabular}
\end{table}

\noindent The thresholds are structurally analogous to the
Herfindahl--Hirschman Index (HHI) tiers used in antitrust analysis
(see~\cite{doj_hhi_2010}), adapted to the ISI's unit-interval scale.
The threshold values are defined in methodology v1.0 and documented
in Paper~1.

\subsection{Rounding and precision}

All axis-level and composite scores are reported to eight decimal
places, matching the precision of the authoritative ISI snapshot.
Population-level statistics (mean, standard deviation, percentiles)
are computed with full floating-point precision and reported to eight
decimal places unless otherwise noted.  The population standard
deviation ($\sigma$, with $N$ denominator) is used throughout, as the
27~EU member states constitute the complete target population.
